% Source: https://artofproblemsolving.com/community/c3h2308030_rectangle_q
% Stable identity: aops-topic-2308030-q01
Let $ABCD$ be a rectangle with diagonal $AC$, and let $\alpha=\angle BAC$, so $\angle CAD=90^\circ-\alpha$.

\begin{asy}
usepackage("amsmath");
size(5.5cm);
defaultpen(fontsize(10pt));
pair A=(0,0), B=(5,0), C=(5,3), D=(0,3);
draw(A--B--C--D--cycle,linewidth(1));
draw(A--C,blue);
draw((4.7,0)--(4.7,0.3)--(5,0.3));
draw(arc(A,0.9,0,degrees(C)),red); draw(arc(A,1.1,degrees(C),90),heavygreen);
label("$\alpha$",1.2*dir(degrees(C)/2),red+fontsize(9pt));
label("$90^\circ-\alpha$",1.15*dir(65),NE,heavygreen+fontsize(9pt));
label("$A$",A,SW); label("$B$",B,SE); label("$C$",C,NE); label("$D$",D,NW);
\end{asy}

In the right triangle $ABC$ (right angle at $B$), $\angle BCA=90^\circ-\alpha$.

\textbf{If $AC$ bisects $\angle A$}, then $\alpha=90^\circ-\alpha$, so $\alpha=45^\circ$. Then triangle $ABC$ has two angles of $45^\circ$, so it is isosceles with $AB=BC$, and the rectangle is a square.

\textbf{If $ABCD$ is a square}, then $AB=BC$, so triangle $ABC$ is isosceles and $\angle BAC=\angle BCA$. These two angles add up to $90^\circ$, so each is $45^\circ$. Hence $\angle BAC=\angle CAD=45^\circ$, and in the same way the diagonal bisects the angle at $C$.

So a diagonal bisects the angles at its endpoints exactly for squares. For a rectangle with $AB>BC$, the larger side is opposite the larger angle in triangle $ABC$, so $\angle BCA>\angle BAC$, that is, $\alpha<45^\circ$: the diagonal is closer to the longer side. $\blacksquare$
